\section{Cohort Methods}

The cohort drew on the records of 310 consecutive patients aged 65 years or older. Mobilization time was taken from nursing notes. Mobilization within 24 hours of surgery counted as early and any later mobilization as usual.

Records were incomplete for 18 patients, who were excluded. Their exclusion may have biased the cohort toward patients with complete nursing notes.

The ethics committee approved the use of the nursing notes after a hassle that lasted four months.

The primary outcome matched the trial outcome. Each patient's readmissions within 30 days of discharge came from hospital records.

The analysis adjusted for age and the Charlson comorbidity index, a weighted count of chronic diseases.

\section{Cohort Results}

Readmission occurred in 22 percent of the 140 patients mobilized early and in 31 percent of the 170 patients with usual mobilization. After adjustment, the odds ratio for readmission was 0.68 (95% CI [0.47, 0.98]), against a crude odds ratio of 0.63. Age and comorbidity probably explained little of the association.

Three cases of deep vein thrombosis were recorded in the cohort, a rate of 1 percent. All three occurred in patients with usual mobilization. Usual mobilization may have added to the risk of thrombosis.

\section{Discussion}

The study adds randomized and observational evidence to the two small trials of early mobilization (Okoye, 2022).

The trial and the cohort agree in direction, which suggests that early mobilization reduces readmission. The trial is randomized; the cohort is observational. Agreement between the two designs strengthens the finding.

Both studies come from one hospital, so the effect may differ elsewhere.

Healthier patients, who could be mobilized earlier, may account for part of the cohort association. Adjustment for age and comorbidity changed the odds ratio only from 0.63 to 0.68, which limits that explanation.

The effect shows in readmission, not in length of stay.

Hospitals that wish to test early mobilization could schedule physiotherapy on the day of surgery.

\section{Conclusion}

Early mobilization within 24 hours may reduce readmission after hip fracture surgery in older adults. Future trials should test it in at least three hospitals.
